\documentclass[10pt,a4paper]{amsart}
\usepackage[margin=19mm]{geometry}
\usepackage{amsmath,amssymb,mathtools}
\usepackage[scr=rsfs,cal=euler]{mathalfa}
\usepackage{xfrac}
\usepackage[unicode,hidelinks,bookmarks=false]{hyperref}
\newcommand{\ie}{i.e.\ }
\newcommand{\cf}{cf.\ }
\newcommand{\resp}{respectively\ }
\newcommand{\Subsection}{\subsection}
\providecommand{\nobreakdash}{\nobreak\hbox{-}}
\setlength{\emergencystretch}{2em}
\multlinegap0pt
\usepackage{bm}
\usepackage{bbm}
\usepackage{dsfont}
\usepackage{xspace}
\usepackage{cancel}
\usepackage{mathtools}
\usepackage{amsthm}
\makeatletter
\@ifundefined{theorem}{\newtheorem{theorem}{Theorem}}{}
\@ifundefined{corollary}{\newtheorem{corollary}[theorem]{Corollary}}{}
\@ifundefined{lemma}{\newtheorem{lemma}[theorem]{Lemma}}{}
\@ifundefined{proposition}{\newtheorem{proposition}[theorem]{Proposition}}{}
\makeatother
\title[Rings Whose Non-Units are Square-Nil Clean]{Rings Whose Non-Units are Square-Nil Clean}
\author{Mina Doostalizadeh, Ahmad Moussavi,, Peter Danchev}
\date{Source: arXiv:2508.01286v1 (CC BY 4.0)}
\begin{document}
\maketitle

\noindent\emph{Source and licence.} Rings Whose Non-Units are Square-Nil Clean. Version arXiv:2508.01286v1 (CC BY 4.0), distributed under the Creative Commons Attribution 4.0 International licence, as recorded in the arXiv licence metadata for this exact version and shown on the abstract page https://arxiv.org/abs/2508.01286v1. Full rights notice: licenses/arxiv-2508.01286-CC-BY-4.0.txt.\par\smallskip

\noindent\textit{Editorial scope.} Counted unit: the Proposition whose source label is \texttt{2.9} in the article (source0.tex lines 369--379, source label \texttt{2.9}), delivered together with the article's own complete original proof of it (source0.tex lines 381--398, one \texttt{proof} environment of 18 lines). No mathematics is written, repaired or re-derived: statement and proof are the article's own text line for line. Required context retained uncounted: 2.4 (lines 276--278); 10 (lines 351--361). Every definition, display and label of those lines is retained unchanged. Omitted from this package and not redistributed: 1--275, 279--350, 362--368, 380--380, 399--919. Nothing inside a retained excerpt is omitted or altered: every retained body is delivered exactly as the article's lines are written, so no omission receipt of any kind accompanies this package. Citations: the retained excerpts carry no citation command at all, so no bibliography entry of the article is transcribed and no reference is redistributed. Reference adaptation: none was needed and none was made; every reference inside the delivered unit resolves inside it, so no unresolved reference or citation is produced. Printed slips kept literal: no printed word, symbol, hypothesis or number of the counted statement or of its proof was altered. Layout and class: the article's own class is \texttt{amsart} with packages including amssymb, amsmath, amscd, amsthm, comment, xcolor, ulem, tikz, float, graphicx, circuitikz, hyperref; the wrapper is the lane's pinned \texttt{amsart} editorial wrapper at 10pt on A4, which restates the theorem-like environments the retained text needs and adds the article's own macro definitions that the retained text needs (none), with extra wrapper packages (bm, bbm, dsfont, xspace, cancel, mathtools). The delivered numbering is the unit's own. Assets: no retained span contains a figure environment, a graphics-inclusion command, a PDF-page inclusion, a table or a TikZ picture; no image file is copied into this package, the package declares no asset dependency and no image file is redistributed with it. Licence: version arXiv:2508.01286v1 of this article is distributed under the Creative Commons Attribution 4.0 International licence, as recorded in the arXiv licence metadata for this exact version and shown on the abstract page https://arxiv.org/abs/2508.01286v1. A full rights notice is delivered at licenses/arxiv-2508.01286-CC-BY-4.0.txt. Characters: every retained excerpt is pure ASCII, so the delivered engine, XeLaTeX with its default Latin Modern text font, reports no missing character.
\begin{lemma}\label{2.4}
Let $R_{i}$ be a ring for all $i\in I$. The direct product $\prod_{i=1}^{n} R_i$ is strongly NUS-nil clean if, and only if, each direct component $R_{i}$ is strongly square-nil clean.
\end{lemma}

\begin{corollary}\label{10}
Let $I$ be an ideal of a ring $R$. Then, the following are equivalent:
\begin{enumerate}
\item
$R/I$ is strongly NUS-nil clean.
\item
$R/I^n$	is strongly NUS-nil clean for all $n \in \mathbb{N}$.
\item
$R/I^n$ is strongly NUS-nil clean for some $n \in \mathbb{N}$.
\end{enumerate}	
\end{corollary}

\begin{proposition}\label{2.9}
Let $R$ be a ring. Then, the following are equivalent:
\begin{enumerate}
\item
$R$ is strongly square nil clean.
\item
${\rm T}_{n}(R)$ is strongly NUS-nil clean for all $n \in \mathbb{N}$.
\item
${\rm T}_n(R)$ is strongly NUS-nil clean for some $n \geq 2$.
\end{enumerate}
\end{proposition}

\begin{proof}
(ii) $\Rightarrow$ (iii). This is trivial.\\
(iii) $\Rightarrow$ (i). Let $a\in R$. Then, it must be that
$$A=\begin{pmatrix}
	a & 0 \\
	0& 0 \\
\end{pmatrix}\in {\rm T}_2(R).$$ It is evident that $A$ is non-invertible in ${\rm T}_2(R)$. By hypothesis, we can find a square idempotent
$E=\begin{pmatrix}
	e_{11} &e_{12} \\
	0& e_{22}  \\
\end{pmatrix}$ and a nilpotent
$Q=\begin{pmatrix}
	q_{11} &q_{12} \\
	0& q_{22}  \\
\end{pmatrix}$ such that
$A=E+Q$ and $EQ=QE$. It now follows by plain inspection that $a=e_{11}+q_{11}$ and $e_{11}q_{11}=q_{11}e_{11}$. Hence, $e_{11}^4=e_{11}^2$, and $q_{11}$ is a nilpotent in $R$. Thus, $a$ has strongly NUS-nil clean decomposition. Therefore, point (i) is valid.\\
(i) $\Rightarrow$ (ii). It suffices with Lemma \ref{10} at hand to prove only that the quotient $T_n(R)/J(T_n(R))$ is strongly NUS-nil clean. To this target, observe that $$T_n(R)/J(T_n(R))\cong \sqcap_{i=1}^nR/J(R).$$ So, we need just to show that $\sqcap_{i=1}^nR/J(R)$ is strongly NUS-nil clean. In fact, Lemma \ref{2.4} is a guarantor that $\sqcap_{i=1}^nR/J(R)$ is strongly NUS-nil clean if, and only if, $R/J(R)$ is strongly square-nil clean. It is now easy to see that, if $R$ is strongly square-nil clean, then $R/J(R)$ is strongly-square nil clean, as asked for.
\end{proof}

\end{document}
